% Safe Fruit -- PTA Programming Assignment Report
% Final program: v8_cycle_tree_dp.c (C99)

\documentclass[12pt,a4paper]{report}
\usepackage[utf8]{inputenc}
\usepackage[T1]{fontenc}
\usepackage{lmodern}
\usepackage{geometry}
\usepackage{amsmath,amssymb,amsthm}
\usepackage{graphicx}
\usepackage{booktabs}
\usepackage{array}
\usepackage{float}
\usepackage{caption}
\usepackage{titlesec}
\usepackage{enumitem}
\usepackage{xcolor}
\usepackage{algorithm}
\usepackage{algorithmicx}
\usepackage{algpseudocode}
\usepackage{indentfirst}
\usepackage{hyperref}

\geometry{left=2.5cm,right=2.5cm,top=2.5cm,bottom=2.5cm}
\setlength{\parindent}{2em}
\setlength{\parskip}{1pt}

\titleformat{\chapter}
  {\normalfont\LARGE\bfseries}
  {\chaptername\ \thechapter.\quad}
  {0pt}
  {}
\titlespacing*{\chapter}{0pt}{0pt}{20pt}

\titleformat{\section}
  {\normalfont\large\bfseries}
  {\thesection}
  {1em}
  {}
\titlespacing*{\section}{0pt}{12pt}{6pt}

\titleformat{\subsection}
  {\normalfont\normalsize\bfseries}
  {\thesubsection}
  {1em}
  {}
\titlespacing*{\subsection}{0pt}{10pt}{4pt}

\setlist[itemize]{left=0pt,topsep=4pt,partopsep=0pt,itemsep=2pt}
\setlist[enumerate]{left=0pt,topsep=4pt,partopsep=0pt,itemsep=2pt}

\hypersetup{colorlinks=true,linkcolor=blue,urlcolor=blue,citecolor=blue}
\newcommand{\code}[1]{\texttt{\detokenize{#1}}}
\newcolumntype{L}[1]{>{\raggedright\arraybackslash}p{#1}}
\newtheorem{theorem}{Theorem}
\newtheorem{lemma}{Lemma}
\setcounter{tocdepth}{1}

\begin{document}

\begin{center}
{\Huge\bfseries Safe Fruit:\\[0.35cm]
 Exact Branch-and-Bound\\[0.35cm]
 with Forest DP\par}
\vspace{1.3cm}
\Large Course: Data Structures and Algorithms\\[.35cm]
\Large Project: Safe Fruit\\[.35cm]
\Large Implementation: C99\\[.35cm]
\Large Completion Date: October 30, 2026\\[1.2cm]
\Large You Longbin \quad 3250106392\\[.4cm]
\Large Huang Yirui \quad 3250104426\\[.4cm]
\Large Chen Zhikun \quad 3250100883\\[1.2cm]
\Large Generated with \LaTeX
\end{center}
\newpage
{\small
\tableofcontents
}
\newpage

\chapter{Introduction}

\section{Description}

The input describes a basket of fruits and a list of warnings.  Each warning
forbids one pair of fruits from being eaten together.  We must choose a subset
of the basket such that no warned pair is chosen.  Among all such subsets, the
program must first maximize the number of chosen fruits and then minimize their
total price.  The problem guarantees that the optimal answer is unique.

The input format is:
\begin{itemize}
\item the first line contains $N$ warnings and $M$ basket fruits;
\item the next $N$ lines each contain a forbidden pair of fruit IDs;
\item the next $M$ lines each contain a fruit ID and its price.
\end{itemize}
Here $N,M\le 100$, every price is between 1 and 1000, and fruit IDs are printed
with three digits.  The output has exactly three lines: the number of chosen
fruits, their IDs in increasing order separated by exactly one space, and the
total price.  There must be no trailing space on the ID line.

Let each basket fruit be a vertex and each warning between two basket fruits be
an undirected edge.  A feasible choice is an independent set of this graph.
The task is therefore a \emph{lexicographic maximum independent set} problem:
maximize the cardinality first, then minimize the sum of vertex prices.

The lexicographic order is important.  A greedy rule that picks the cheapest
fruit first can fail.  For example, on the path $1-2-3$ with prices
$p_1=100$, $p_2=1$, $p_3=100$, the cheapest fruit is 2, but choosing it gives
only one fruit; the optimal choice is $\{1,3\}$ with two fruits and total price
200.  The program must therefore optimize the count before the price, not the
price alone.

General maximum independent set is NP-hard, so no polynomial-time exact
algorithm is expected.  The assignment limits $M$ to 100, which makes an exact
branch-and-bound search with strong pruning practical.

A warning may mention a fruit that is not in the basket.  Such a warning does
not constrain the selection, and the program ignores it defensively.  The PTA
sample does not actually contain one: its first 16 lines are warnings, and the
lines \code{020 99} and \code{019 99} that follow are fruit lines, meaning that
fruits 020 and 019 each cost 99.  The special tests in
Chapter~\ref{ch:testing} include outside warnings explicitly.

The sample has 16 warnings and 20 fruits.  The warnings form three nontrivial
components: the edge $001-002$; the 6-cycle
$003-004-005-006-007-008-003$; and the tree on $009$--$018$, where 009 is
adjacent to 010, 011, 012, 013, and 010-014, 011-015, 012-016, 012-017,
013-018 are the remaining edges.  Fruits 019 and 020 have no warnings and cost
99 each.  The sample answer is
\begin{verbatim}
002 004 006 008 009 014 015 016 017 018 019 020
\end{verbatim}
with total price 239.  The choices 002, 004, 006, and 008 come from the first
two components; 009 and 014--018 come from the tree; 019 and 020 are isolated
and are selected because they increase the count.  The price check is
$1+4+5+1+10+5+6+3+2+4+99+99=239$.  This is a useful sanity check for the score
encoding and for the isolated-fruit rule.

The final program should be exact, implement the count-first, price-second
objective, print the required format, and run comfortably for $M\le100$.  The
submitted program is \code{v8_cycle_tree_dp.c}.  Earlier versions are kept in
the repository because they were used to validate each optimization during
development; they are not alternative answers.

\section{Overview of the Implementation}

The program is organized around four components.
\begin{enumerate}
\item \textbf{Graph construction and PTA input.}  The program reads $N$ and $M$,
      maps the three-digit IDs to internal indices, and builds both a bitset
      adjacency and adjacency lists.  Isolated fruits are added to the answer
      directly.  A warning with an endpoint outside the basket is ignored; this
      is defensive, because the sample has no such warning.
\item \textbf{Score encoding and search.}  Each fruit receives the weight
      $B-p_v$, so one integer captures both objectives.  The search uses two
      64-bit masks for the candidate and selected sets, a matching upper bound,
      and a dynamic maximum-degree pivot.
\item \textbf{Structural decomposition.}  Original connected components are
      solved independently.  A tree component uses tree DP; a cyclic component
      uses branch and bound, with residual-forest DP when the candidate graph
      becomes acyclic.  Version 8 chooses when to perform the forest check.
\item \textbf{Output.}  The selected IDs are sorted in increasing order and
      printed with \texttt{\%03d}; a space is printed before every ID except
      the first, so the line has no trailing space.  The second line is followed
      by the total price on the third line.
\end{enumerate}
The count on the first line is the number of selected IDs, and the price is
their sum.  The input is assumed to be a simple conflict graph, as the
assignment states that there are no duplicated tips; self-loops are not
specified.

The final program is compiled with \code{gcc -O2 -std=c99}; the optional
\code{-DBENCH} macro is used only for local timing and is not needed by the
judge.  Chapter~\ref{ch:algorithm} specifies the algorithm and its correctness;
Chapter~\ref{ch:testing} is the test report; Chapter~\ref{ch:analysis} gives the
analysis and comments; the appendix gives the source location.

\chapter{Algorithm Specification}\label{ch:algorithm}

\section{Encoding the Two Objectives}

Let $p_v$ be the price of fruit $v$.  Define
\[
B=100001,\qquad w_v=B-p_v,\qquad W(S)=B|S|-\sum_{v\in S}p_v .
\]
The maximum possible total price is $100\times1000=100000<B$, and every
$w_v$ is positive.

\begin{lemma}
For any two feasible sets $A$ and $C$, $W(A)>W(C)$ if and only if
\[
\Bigl(|A|,-\sum_{v\in A}p_v\Bigr)
\quad\text{is lexicographically greater than}\quad
\Bigl(|C|,-\sum_{v\in C}p_v\Bigr).
\]
\end{lemma}

\begin{proof}
If $|A|>|C|$, then
\[
W(A)-W(C)=B(|A|-|C|)-\Bigl(\sum_{A}p-\sum_{C}p\Bigr)
\ge B-1000|A|\ge B-100000=1>0 .
\]
If $|A|=|C|$, then $W(A)-W(C)=\sum_{C}p-\sum_{A}p$, so a larger score means a
smaller total price.
\end{proof}

Thus maximizing the single integer $W$ is exactly equivalent to the required
lexicographic objective.  The largest possible score is at most
$100B=10{,}000{,}100$, which fits safely in a signed 32-bit integer.  The score
is additive over disjoint components and over selected vertices, so the same
encoding can be used by component decomposition and tree DP.

\section{Branch and Bound}

A search state consists of the selected set $S$, its score $s$, and the
candidate set $C$ of fruits that are compatible with every vertex of $S$.
For a pivot $v\in C$, every feasible completion belongs to exactly one of two
cases:
\[
\begin{array}{ll}
\text{take }v:& (S\cup\{v\},\ s+w_v,\ C\setminus(\{v\}\cup N(v))),\\
\text{skip }v:& (S,\ s,\ C\setminus\{v\}).
\end{array}
\]
The take branch is searched first.  This is a heuristic for finding a good
incumbent early; it is not required for correctness.

The search is pruned with an upper bound $\operatorname{UB}(C)$ on the score
that can still be added from $C$.  If $s+\operatorname{UB}(C)\le W_{\text{best}}$,
the state cannot strictly improve the incumbent and is discarded.  Equality is
safe because the encoded score already includes the price tie-break and any
optimal set is acceptable.

The pivot is chosen dynamically as a vertex with the largest number of
neighbors still in $C$.  Its take branch removes the pivot and all its remaining
neighbors, which often reduces the search tree.  Its skip branch removes only
the pivot, so the rule is a heuristic rather than a guarantee.  A static
degree order is cheaper per node but becomes less informative as candidates
disappear; the final program uses the dynamic rule.

\section{Matching Upper Bound}

Summing all candidate weights is a valid but loose upper bound.  A tighter
bound is obtained from a greedy matching $Q$ in the induced conflict graph
$G[C]$.  Let $U$ be the unmatched candidates.  Then
\[
\operatorname{UB}(C)=\sum_{(u,v)\in Q}\max(w_u,w_v)+\sum_{u\in U}w_u .
\]

\begin{lemma}
$\operatorname{UB}(C)$ is an upper bound on the score of any independent subset
of $C$.
\end{lemma}

\begin{proof}
The pairs of $Q$ are disjoint.  An independent set can contain at most one
endpoint of each matched pair, so the contribution of a pair is at most the
larger of its two weights.  An unmatched vertex contributes at most its own
weight.  Conflicts between different pairs are ignored, which can only make the
bound larger.  Hence the sum is an upper bound.
\end{proof}

No maximum matching is required for correctness.  The implementation builds the
matching greedily, scanning vertices in decreasing initial degree order.  A
different order changes the tightness of the bound but not its validity.

\section{Component Decomposition}

If the conflict graph has connected components $G_1,\ldots,G_k$, then there are
no edges between different components, and
\[
\operatorname{OPT}(G)=\sum_{i=1}^{k}\operatorname{OPT}(G_i)
\]
for the additive score.  The program solves each original component separately
and combines the selected masks.  This avoids exploring the Cartesian product
of independent local choices, which is especially valuable when several
components have many tied local optima.

The program does not re-split arbitrary cyclic candidate graphs at every
recursive node.  However, the forest routine described below can solve several
residual trees together inside one component.

\section{Tree DP}

A connected simple graph with $n$ vertices and $n-1$ edges is a tree.  For a
rooted tree, let $T(v)$ be the best score of the subtree rooted at $v$ when $v$
is selected, and $F(v)$ the best score when $v$ is excluded.  Then
\[
T(v)=w_v+\sum_{u\text{ child of }v}F(u),\qquad
F(v)=\sum_{u\text{ child of }v}\max\{T(u),F(u)\}.
\]
The values are computed in reverse traversal order and a choice is reconstructed
from the root downward.  A selected parent forces every child to be excluded;
otherwise the better child state is chosen.  The optimum of the tree is
$\max\{T(\text{root}),F(\text{root})\}$.

This DP is exact by induction on the height of the subtree.  The final program
applies it to every original component that is a tree, so paths, stars, and
trees are solved without branching.

\section{Residual Forest DP}

Branching removes vertices, so a cyclic component can become a forest.  During
the search, the program tests the current candidate graph $G[C]$.  A traversal
checks acyclicity and establishes parent pointers.  If a cycle remains, ordinary
search continues.  If $G[C]$ is a forest, the same take/skip recurrence solves
every residual tree, the resulting score is added to $s$, and that search branch
ends immediately.

This substitution is exact.  Every residual candidate is compatible with $S$,
there are no edges between different residual trees, and the DP is exact on each
tree.  The algorithm never approximates a cyclic graph by a tree; it explicitly
verifies acyclicity before using the recurrence.

\section{When to Check: Versions 7 and 8}

A forest test costs $O(M+E)$ even when it fails.  Testing every search node can
waste traversal work while many cycles remain; testing only very late can miss
an inexpensive DP completion.  Two schedules were developed:
\begin{itemize}
\item Version 7 uses fixed thresholds.  When the original component has cycle
      rank $r=e-n+1<3$, it first checks after at least five candidates have
      disappeared; when $r\ge3$, it first checks only when at most 15
      candidates remain.  After a failed check it waits for five more
      disappearances along that branch.
\item Version 8 sets the first threshold from the component's cycle rank:
      \[
      q=\min\{n,\lceil r/2\rceil\},
      \]
      and performs the first test after at least $q$ candidates have
      disappeared.  After a failed test, a descendant is tested only when its
      incoming branch removed at least two candidates.  A take branch usually
      removes a pivot and several neighbors; a skip branch removes one vertex
      and therefore does not immediately trigger another check.
\end{itemize}

Versions 7 and 8 use the \emph{same} forest-DP substitution and the same
correctness argument.  They differ only in when the test is attempted.  In
other words, v8 is a parameter choice inside v7's mechanism, not a new
algorithm.

The schedule is heuristic.  The cycle rank counts independent cycles, not the
number of vertices that must be deleted: removing one hub can destroy many
cycles, while deleting unrelated leaves may destroy none.  The rank is computed
once per original component and is not recomputed at every state.  Delaying or
omitting a test cannot lose an optimum; it only leaves more work to exact branch
and bound.

A concrete limitation of the formula is that $q=n$ when $r\ge 2n-1$.  Then the
first test is deferred until $C=\varnothing$, so a dense high-rank component gets
no early forest detection at all.  For example, a 20-vertex component with 100
edges has $r=81$ and $q=20$.  This is one reason the schedule should be read as
a heuristic parameter, not an optimal rule.  On some graphs v8 can be slower
than v7, because the extra checks can cost more than the search they save.
Correctness is unaffected, because a delayed test only leaves the work to the
exact search.

\section{Final Algorithm}

Algorithm~\ref{alg:search} shows the final search for one cyclic component.
The component size $n$ and the first threshold $q$ are fixed before the call.
Isolated fruits are added to the answer directly and are not part of $C$.

\begin{algorithm}[H]
\caption{Structural branch and bound (v8, one cyclic component)}
\label{alg:search}
\begin{algorithmic}[1]
\Procedure{Search}{$C,s,S,\mathit{started},\Delta$}
\If{$s+\operatorname{UB}(C)\le W_{\mathrm{best}}$}\State\Return\EndIf
\If{$(\neg\mathit{started}\land n-|C|\ge q)\lor(\mathit{started}\land\Delta\ge2)$}
  \If{$G[C]$ is a forest}
    \State Solve all residual trees by DP and update the incumbent
    \State\Return
  \EndIf
  \State $\mathit{started}\gets\mathrm{true}$
\EndIf
\If{$C=\varnothing$}\State Update incumbent with $(s,S)$; \Return\EndIf
\State $v\gets$ maximum-degree vertex in $G[C]$
\State $C'\gets C\setminus(\{v\}\cup N(v))$
\State \Call{Search}{$C',s+w_v,S\cup\{v\},\mathit{started},|C|-|C'|$}
\State \Call{Search}{$C\setminus\{v\},s,S,\mathit{started},1$}
\EndProcedure
\end{algorithmic}
\end{algorithm}

\section{Correctness}

\begin{theorem}
For every valid PTA input, the program outputs a feasible set of maximum
cardinality; among all such sets it outputs one of minimum total price.
\end{theorem}

\begin{proof}
By Lemma~1, maximizing $W$ is equivalent to the required lexicographic
objective.  For a pivot $v$, every independent set either contains $v$ or does
not; the take and skip branches therefore partition all feasible completions.
The matching bound is a valid upper bound by Lemma~2, so a state is discarded
only when it cannot strictly improve the incumbent.  Component decomposition is
valid because the score is additive and there are no cross-component edges.
The tree and forest DPs are exact by induction on subtree height and are invoked
only after the acyclicity check.  Therefore the search never discards an
optimal set, and the reconstructed mask is feasible and optimal.  The final
sorting by ID only changes the output order, not the selected set.
\end{proof}

\section{Complexity}

Let $K$ be the number of visited search states, $M$ the basket size, and $E$
the number of represented conflict edges.  The matching bound, the dynamic
pivot, and a triggered forest check each cost $O(M+E)$ in a conservative bound.
The total time is therefore $O(K(M+E))$, with $K=O(2^M)$ in the worst case.

The adjacency matrix uses $O(M^2)$ space and the recursion $O(M)$ stack space;
with $M\le100$ the dense array is small and two 64-bit words cover all
candidates.  The worst case remains exponential, but the pruning and structural
routines make the PTA limits comfortable.

\chapter{Testing Results}\label{ch:testing}
\section{Test Cases and Input Data}

The cases cover the PTA sample, boundary sizes, the main graph shapes, the
lexicographic objective, outside-basket warnings, and output formatting.
Table~\ref{tab:tests} lists them.  The two large cases are constructed as
follows:
\begin{itemize}
\item \textbf{K14}: all $\binom{14}{2}=91$ pairs are warnings; prices are
      1, 2, \ldots, 14.  The optimum is the single cheapest fruit.
\item \textbf{Path100}: warnings form the path
      $001-002-\cdots-099-100$, plus the outside warning \code{001 999};
      odd IDs have price 1 and even IDs have price 100.  The unique optimum is
      the 50 odd IDs, with total price 50.
\end{itemize}
The other cases are small enough to be read directly from their descriptions.

\begin{table}[H]\centering\scriptsize
\setlength{\tabcolsep}{4pt}
\renewcommand{\arraystretch}{1.15}
\caption{Test cases and results.  ``same'' means the actual output was
byte-for-byte identical to the expected output.}\label{tab:tests}
\begin{tabular}{@{}c L{2.9cm} c L{4.1cm} L{4.1cm} c@{}}
\toprule
No. & Purpose & $N/M$ & Expected (count / IDs / price) & Actual & Result\\
\midrule
1 & PTA sample; isolated fruits 019/020 & 16/20 &
\texttt{12 / 002 004 006 008 009 014 015 016 017 018 019 020 / 239} & same & Pass\\
2 & Single fruit; outside tip & 1/1 &
\texttt{1 / 001 / 5} & same & Pass\\
3 & All isolated; every fruit selected & 3/5 &
\texttt{5 / 001 002 003 004 005 / 150} & same & Pass\\
4 & Complete $K_4$; count 1, cheapest & 6/4 &
\texttt{1 / 001 / 10} & same & Pass\\
5 & Path with cheap middle; count first & 2/3 &
\texttt{2 / 001 003 / 200} & same & Pass\\
6 & Path; same count, lower price & 3/4 &
\texttt{2 / 001 004 / 2} & same & Pass\\
7 & Two disjoint edges; components & 2/4 &
\texttt{2 / 001 003 / 4} & same & Pass\\
8 & 4-cycle; same count, lower price & 4/4 &
\texttt{2 / 001 003 / 4} & same & Pass\\
9 & Star; leaves beat the centre & 3/4 &
\texttt{3 / 002 003 004 / 60} & same & Pass\\
10 & Leading zeros 010 and 100 & 1/2 &
\texttt{2 / 010 100 / 12} & same & Pass\\
11 & $K_{14}$; 91 tips, dense & 91/14 &
\texttt{1 / 001 / 1} & same & Pass\\
12 & Path100; max $N=M=100$ & 100/100 &
\texttt{50 / 001 003 ... 097 099 / 50} & same & Pass\\
\bottomrule
\end{tabular}
\end{table}

\section{Measurement Method}

The tests were run on Windows with gcc 15.2.0 (MinGW-w64), using
\code{gcc -O2 -std=c99}.  Each case was run as a separate process with the test
input on standard input.  The output was compared with the expected output, and
a checker verified:
\begin{itemize}
\item the first line equals the number of IDs on the second line;
\item the IDs are sorted, unique, and present in the basket;
\item no forbidden pair is selected;
\item the third line equals the sum of the selected prices;
\item the count and price equal the expected optimum.
\end{itemize}

\section{Sample Run}

The exact output for the PTA sample is shown below.  It matches the required
format: three lines, one space between IDs, no trailing space.
\begin{verbatim}
12
002 004 006 008 009 014 015 016 017 018 019 020
239
\end{verbatim}

\section{Additional Automated Tests}

Two repository scripts provide broader coverage:
\begin{itemize}
\item \code{test_random.py} generates 225 random small cases, computes the
      optimum by enumerating all subsets, and checks every version.  The final
      program and all earlier versions pass.
\item \code{test_special.py} checks five special cases: objective priority,
      multiple optima, dirty input (outside tips and a duplicate tip), multiple
      components, and a maximum-size case with 50 disjoint edges.  All versions
      pass.  The multiple-optima case is stricter than the PTA guarantee, which
      says the answer is unique; it verifies that any optimal set is accepted
      by the checker.
\end{itemize}
In total, the final program passed the 12 hand-written cases, the 225 random
cases, and the 5 special cases.  The sample output and the formatting checks
also passed.

\section{Performance Observation}

Table~\ref{tab:perf} gives an illustrative single-run measurement of the
internal solve time for the final program and the v2 baseline.  It is not a
benchmark suite; it only shows the effect of the structural optimizations.
\begin{table}[H]\centering\small
\caption{Illustrative internal solve time on the test machine (one run).}
\label{tab:perf}
\begin{tabular}{@{}lrr@{}}\toprule
Case & v2 baseline (ms) & v8 final (ms)\\\midrule
Path100 & 6.44 & 0.0016\\
Cycle100 & 10.93 & 0.0035\\
$K_{14}$ & 0.0024 & 0.0034\\
\bottomrule\end{tabular}
\end{table}
The path and cycle cases show the benefit of tree and residual-forest DP: the
baseline still branches, while the final program finishes the acyclic remainder
directly.  The tiny $K_{14}$ case shows the opposite trade-off: the graph is so
easy that the extra forest checks can cost slightly more than the search they
save.  This is a performance detail, not a correctness issue, and it is the same
reason the v8 schedule should not be described as universally better than v7.

\chapter{Analysis and Comments}\label{ch:analysis}

\section{Why the Algorithm Is Exact}

The proof has four independent parts.  The score encoding turns the two
objectives into one integer.  The take/skip branches partition all independent
sets.  The matching bound is a valid upper bound, so pruning never removes a
state that could improve the incumbent.  Component decomposition and tree or
forest DP are exact on their subproblems.  The forest check is separated from
the DP: the check may be delayed, but when it succeeds the DP is exact, and when
it fails the ordinary search continues.

This separation is the main design idea.  A heuristic trigger is allowed to
affect speed only; it cannot affect the answer.

\section{Why the Optimizations Matter}

The matching bound discards score combinations that cannot beat the incumbent.
Component decomposition avoids the Cartesian product of independent local
choices.  The dynamic pivot removes a high-degree vertex and many neighbors in
the take branch.  Tree DP replaces enumeration on a path, star, or tree, and
residual-forest DP extends the same idea to the acyclic remainder created by
branching.

On the 100-vertex path and cycle, the final program solves the problem in a few
microseconds of internal time, while the baseline branches for milliseconds.
The improvement is not a constant-factor tweak; it changes the work required
for a whole tractable subproblem.

\section{PTA-Specific Pitfalls}

\begin{enumerate}
\item \textbf{Outside warnings (defensive).}  A warning whose endpoint is not
      in the basket is ignored.  The PTA sample does not exercise this rule
      (the lines \code{020 99} and \code{019 99} are fruit lines), but the
      special tests do.  Ignoring such a warning is safe, because a fruit that
      is not in the basket cannot be selected.
\item \textbf{Count before price.}  A cheap fruit may block two expensive
      fruits.  The score $B-p_v$ makes the count dominant; do not sort by price
      or use a price-only greedy rule.
\item \textbf{Unique answer.}  The problem guarantees a unique optimum, but the
      program still implements the price tie-break.  It does not need to
      enumerate all optimal sets.
\item \textbf{Three-digit IDs.}  Print with \texttt{\%03d}; sort numerically, not
      lexicographically.  The leading-zero case \code{010}/\code{100} checks
      this.
\item \textbf{Output format.}  Exactly three lines; IDs in increasing order;
      exactly one space between IDs; no trailing space.  The count must equal
      the number of IDs, and the price must equal their sum.
\item \textbf{Maximum size.}  $M=100$ fits in two 64-bit masks.  The maximum
      total price is 100000, so $B=100001$ is large enough; the maximum score
      fits in a signed 32-bit integer.
\item \textbf{Isolated fruits.}  A fruit with no warning is always selected,
      because its weight is positive.  The program adds isolated fruits
      directly and removes them from the search.
\item \textbf{Input assumptions.}  The statement says there are no duplicate
      tips.  The program assumes a simple graph; it does not specially handle
      self-loops, which the statement does not define.
\end{enumerate}

\section{Limitations}

The worst case is still exponential: maximum independent set is NP-hard, and an
adversarial graph with $M=100$ can make branch and bound explore many states.
The program is exact, so it will eventually finish, but a malicious test could
be slow.  The forest-detection schedule is a heuristic; its parameters were
chosen to work well on the tests, and v8 can be slower than v7 on some graphs,
especially when branching is already cheap.  This is a performance trade-off,
not a correctness trade-off.  The v7 schedule is a simpler alternative if
robustness across graphs is more important than aggregate speed.

\section{Lessons Learned}

\begin{itemize}
\item Start from a correct baseline and add one optimization at a time; test
      after each change.
\item Use brute force on small random cases to validate the objective and the
      tie-break, not just the sample.
\item Keep the exact structural solver separate from the heuristic that decides
      when to call it.  This makes the correctness argument independent of the
      tuning.
\item Output formatting is part of the PTA contract: one extra character can
      fail an otherwise correct solution.
\item Performance measurements are useful, but a single observation on one
      machine should not be over-interpreted.
\end{itemize}

\section{Summary}

The final program solves the PTA Safe Fruit problem exactly.  It encodes the
count-first, price-second objective as one integer, prunes with a matching
bound, decomposes connected components, and uses tree and residual-forest DP to
avoid branching on acyclic structure.  It handles the PTA-specific cases:
defensive handling of outside warnings, leading-zero IDs, maximum size, and the
exact output format.

The program passed the sample, 12 designed cases, 225 random brute-force cases,
and 5 special cases.  Versions 7 and 8 share the same forest-DP mechanism; v8
only changes when the forest check is attempted.  The v8 schedule is faster in
aggregate on the tested cases, but it is a parameter choice, not a universally
better algorithm, and some graphs can regress.  For the PTA assignment, the
final program is exact, well tested, and comfortably fast for $M\le100$.

\chapter*{References}
\addcontentsline{toc}{chapter}{References}
\begin{enumerate}[label={[\arabic*]}]
  \item \textit{Safe Fruit}, PTA programming assignment, Zhejiang University.
  \item T. H. Cormen, C. E. Leiserson, R. L. Rivest, and C. Stein.
        \textit{Introduction to Algorithms}, 4th edition. MIT Press, 2022.
\end{enumerate}

\appendix
\chapter{Source Location}

The final program and the development versions are available at
\begin{center}
\url{https://github.com/blxrua/ads_project3_ads}\quad (commit \code{d2614e2}).
\end{center}
The report source, compiled PDF, and the 12 hand-written test inputs are
included in the repository root and the \code{tests/} directory.
The relevant files are:
\begin{itemize}
\item \code{report_safe_fruit.tex} and \code{report_safe_fruit.pdf}: this report;
\item \code{v8_cycle_tree_dp.c}: the final program;
\item \code{v1_steps123.c} through \code{v7_cycle_tree_dp.c}: development
      history used for regression testing;
\item \code{test_random.py}: 225 random cases checked against brute force;
\item \code{test_special.py}: five special cases, including dirty input and the
      maximum-size case;
\item \code{sample.in}: the PTA sample input.
\end{itemize}

The final program was compiled with
\begin{verbatim}
gcc -O2 -std=c99 -o safe_fruit v8_cycle_tree_dp.c
\end{verbatim}
and run with the input on standard input.  The optional timing build uses
\code{-DBENCH}; it is not needed for the PTA submission.

\end{document}
